\paragraph{The filtered product decomposition.}

Interpolation gives values on the cut. To estimate the cokernel of multiplication we need to express a section through products of short and light sections in successive degrees.

Write \(a=T-s\), and let \(t'_i\) be the chosen short
rational forms of weight \(w_F\) whose restrictions
are a basis on the cut. For a section \(f\) of weight
\(\kappa=w_F+Tw_L\) with the low modes through \(c_1\)
zero, match its jets of order less than \(a\) using
the products
\[
            t'_i L_0^{T-|\alpha|}\prod_j S_j^{\alpha_j},
           \qquad |\alpha|<a .
\]
Divide by \(L_0^T\) locally at every cut point and use
the coordinates \(S_i/L_0\). Let \(\mathcal E\) be the evaluation matrix of the \(t'_i\) on
the entire cut.

If \(v_\alpha\) is the residual Taylor coefficient
on the cut for \(|\alpha|=j<a\), solve
\[
       \mathcal E\,b_\alpha=v_\alpha,\qquad
       g_{j+1}=g_j-\sum_{|\alpha|=j}\sum_i(b_\alpha)_i\,
             t'_i L_0^{T-j}\prod_m S_m^{\alpha_m},
       \qquad g_0=f .
\]
The new terms lie in the \(j\)-th power of the cut ideal, so
\(\mathcal E\) inverts the degree-\(j\) coefficient without
changing any coefficient of smaller degree. This leaves a section \(g\) vanishing
to order at least \(a\) along the whole cut. The goal
is to divide \(g\) by the degree-\(a\) monomials
with coefficients of weight \(w_F+s w_L\), bounded
and cuspidal on the quotient.

Put \(W_L=\sum_l\|w_l\|^2\) and
\(S=\sum_i\|S_i\|^2\) in the normalized light metric.
Take the quotient integral of \(|g|^2\) with the
weight
\[
     W_{\rm div}=e^{-\varphi_\kappa}
                           W_L^{-1/2}S^{-(T-1/2)} .
\]
Near a cut point use the complex dimension \(s\) when integrating
the bound from the order-\(a\) vanishing:
\[
   \int_0^\epsilon r^{\,2a-2(T-1/2)}r^{\,2s-1}\,dr=\epsilon .
\]
On the complement \(W_L\ll S\) bounds the singular factor
by a constant times \(W_L^{-T}\).
The \hyperref[par:slice-cusp-decay]{cuspidal decay estimate}
on complete rank-one periods at total cutoff \(c_1\) is
\(\|f(x)\|_{\kappa}^2\le C e^{-2\pi c_1/\lambda_1(x)^2}\).
The selected cutoff has \(TcP\le c_1/\lambda_1^2\),
where \(P\) is the first-minimum penalty of the light
family.

In detail \(P\le2 f/\lambda_1^2\) and
\(c_1\ge 2fTc\). The lower bound \(W_L\gg e^{-2cP}\) gives uniform control
on the complement of the cut in the cuspidal region:
\[
   \|f\|_\kappa^2 W_L^{-T}
    \ll \exp(2TcP-2\pi c_1/\lambda_1^2)
    \le \exp\bigl(-2(\pi-1)c_1/\lambda_1^2\bigr).
\]
Each correcting product has squared norm \(\ll W_L^T\):
the \(t'_i\) are bounded, and the squared norm of each light
factor is bounded by a constant times \(W_L\). The quotient has finite invariant
volume. For the definite lift compactness and
nonvanishing of the light family give this integrability
with the full cusp space.

\paragraph{Lifting before dividing.}
Let \(k=e^{-\beta R}\), using the exponent \(\beta\) of Lemma~\ref{lem:short-weight-curvature}.
Poincar\'e summation maps the holomorphic functions
of finite norm
\(\int |F|^2 W_{\rm div}D/k\) on the open flag cell to
a subspace of the quotient weighted Hilbert space.
After automorphy factors are absorbed into the norms, the
pointwise Cauchy--Schwarz inequality is
\[
 \Bigl|\sum_\gamma F(\gamma x)\Bigr|^2
 \le \Bigl(\sum_\gamma k(\gamma x)\Bigr)
       \sum_\gamma |F(\gamma x)|^2/k(\gamma x).
\]
Integration on the quotient unfolds the second sum to the
whole-cell integral, and the first sum is uniformly bounded.
The ball bound gives holomorphy and boundedness because
\(W_{\rm div}\gg e^{-\varphi_\kappa}\). The weight-gap
cuspidality criterion at \(\kappa\) places the image in
the finite-dimensional cusp space, so the image is closed
in the quotient weighted norm. We use the same criterion
for bounded holomorphic sections at \(w_F+s w_L\).
The within-block and extreme cross-block comparisons
in the zero-mode calculation apply at both weights.

\begin{proposition}[Analytic lifting and division on the common cut]
\label{prop:lift-and-divide-at-slice}
Let \(f\) be a section of weight \(w_F+Tw_L\) whose low modes
through \(c_1\) vanish, or a section in the full cusp space on the
compact definite lift. Assume the \(t'_i\) of weight \(w_F\) give
a basis on the whole cut, that the \(S_j/L_0\) are coordinates
there, and use the division margin
\(\beta^2u_0+\beta<m_{\rm div}\) of
Lemma~\ref{lem:short-weight-curvature} and the cell estimate of
Lemma~\ref{lem:short-poincare-summation}. Put \(a=T-s\).
Let \(J\) denote the section matching the jets through degree \(a-1\). The difference lies in the \(a\)-th power of the cut ideal and has an expansion
\[
 f-J=\sum_{|\alpha|=a}S^\alpha h_\alpha,\qquad \operatorname{wt}(h_\alpha)=w_F+s w_L,
\]
with every \(h_\alpha\) a bounded cusp form.
\end{proposition}

\begin{proof}
\proofstep{Surjectivity of the automorphic Poincaré sum.}
The next calculation proves surjectivity onto cusp
forms with finite quotient norm for this singular
weight. Suppose \(g_0\) in that space is orthogonal
to the image. Unfolding with the summable kernel gives
\[
          I=\int k|g_0|^2 W_{\rm div}D<\infty
\]
on the cell. In these coordinates write
\(\Psi=-\log W_{\rm div}-\log D\). The formula for
\(W_{\rm div}\), the nonnegative Levi forms of the
two logs of sums of squares, and the division
margin give 
\[
  \mathcal L(\Psi)\ge m_{\rm div}\mathcal L(R)
\]
where the cut functions do not vanish simultaneously.
Work on the pseudoconvex cell after removing the
zero hypersurface of one nonzero cut function.

\proofstep{Transferring the radius to the unfolded cell.}
The radius is pulled back from the bounded matrix
domain with potential \(-\log\det(I-Z^*Z)\) at each
embedding. Its gradient squared in its Levi form is
\(\operatorname{tr}(Z^*Z)\): at diagonal singular
values \(z_i\) the gradient entries are
\(z_i/(1-z_i^2)\) and the corresponding Levi
coefficients are \((1-z_i^2)^{-2}\).
This is at most \(\min(p,q)\) since \(\|Z\|_{\rm op}<1\).
Cauchy--Schwarz across the embeddings and pullback
to the flag cell give

\[
  |\partial R(\xi)|^2\le u_0\mathcal L(R)(\xi)
\]
for every tangent vector \(\xi\).
Apply the weighted \(\bar\partial\) estimate with
potential \(\Psi-\beta R\) to the derivative of
\(kg_0\). Its datum is \(-\beta k g_0\bar\partial R\),
and the inverse Levi cost is bounded using
\((m_{\rm div}-\beta)\mathcal L(R)\).
This gives a holomorphic \(F\) such that
\[
 \int |kg_0-F|^2\frac{W_{\rm div}}{k}D
       \le \frac{\beta^2u_0}{m_{\rm div}-\beta}I .
\]
The Levi form can be semidefinite in the compact flag
directions because the datum vanishes in those directions.
Use exhaustion and the directional estimate to apply the
weighted Hilbert space argument. Since the scalar on the
right is strictly less than one, \(F\) has finite seed
norm. It extends across the removed hypersurface by local
square integrability, using the local positive lower
bound for the metric. It extends to equivariant readings
on the other frame charts in the same way as the
interpolating seed.

Orthogonality to the sum of this \(F\) implies

\[
  \int F\overline{g_0}W_{\rm div}D=0
\]
 by unfolding.
The integrals are absolutely convergent by the seed norm
and \(I<\infty\). Thus the left hand side in the displayed
inequality equals \(I+\int |F|^2 W_{\rm div}D/k\).
It follows that \(I=0\), so \(g_0=0\) by the positive
weight almost everywhere and holomorphy. This shows that
\(g\) has a seed in the specified whole-cell norm.
The stronger inequality 
\[
  \beta^2u_0+\beta<m_{\rm div}
\]
enters here.
\cite{hormander-slice}\label{cor:slice-short-division-1}

\paragraph{Iterating division and returning to the quotient.}

After lifting once, division leaves a section of lower degree. The norm change at each repetition determines the cost on returning to the quotient.

In a local holomorphic frame write
\(S_0'=\sum_i|S_i|^2\), \(W_0'=\sum_l|w_l|^2\).
The seed norm for \(g\) is
\[
  \int |F|^2 (S_0')^{-(T-1/2)} e^{-\Phi},\qquad
     \Phi=\varphi_F-\log D+\tfrac12\log W_0'-\beta R .
\]
The light family has no common zero, and the division
margin makes \(\Phi\) plurisubharmonic. Use Skoda's
\(L^2\) division theorem with \(s\) holomorphic
generators in dimension \(s\). For an exponent
\(H>s\), finite norm with weight
\((S_0')^{-H}e^{-\Phi}\) gives division by the \(S_i\)
with coefficients of finite norm with exponent \(H-1\).
In the theorem's notation \(q=s-1\) and

\[
  H=\alpha q+1
\]
 with \(\alpha>1\).
Here the last input exponent is \(s+1/2\), so all
\(a=T-s\) divisions are allowed. Sum the finitely
many ordered products with the same exponent tuple.
Each resulting seed has the weight with exponent
\(s-1/2\). It extends across the removed zero
hypersurface by local square integrability.

\proofstep{Convergence and descent after division.}
The \(S_i\) and \(w_l\) are bounded in their normalized metric. Consequently each resulting seed \(u\) has finite norm for the regular metric, with the indicated density,
\[
  \int|u|^2 e^{-\varphi_{w_F+s w_L}}\frac{D}{k}<\infty.
\]
Extend each seed across the frame chart divisor
and Poincar\'e sum it in weight \(w_F+s w_L\).
The resulting forms are bounded by the ball estimate
and are cuspidal by the shifted-weight criterion.
Multiplication by each automorphic monomial
\(\prod_i S_i^{\alpha_i}\) commutes with summation
in the corresponding weight. Absolute convergence
therefore yields the division equation for \(g\)
on the whole quotient. The construction also works
in the definite case: \(R=0\), the gradient cost
is zero, and the compact quotient has no cusps.
\cite{skoda-slice}\label{cor:slice-short-division-2}

\end{proof}

\paragraph{Analytic realization of theta sections at the shifted weights.}\label{par:shifted-theta}
Fix one of the two runs and its embedding \(\sigma_0:K\to\C\).
Write \(C(u)\) for the space of rational cusp forms at weight \(u\).
We need three consequences of the arithmetic realization. The analytic realizations of \(C(w_F+k w_L)\) span the full space of complex cusp families for
every \(k\ge0\). At \(\kappa=w_F+T w_L\) the subspace with the
specified low modes zero is spanned by the analytic images of the rational good space. Finally at \(w_F,w_F+s w_L,\kappa\) there are spanning
sets in the prescribed scaled theta lattices with a bound on the mean
log supremum of each member. The last assertion is a size estimate.

\begin{proposition}[Theta spanning on marked components]\label{prop:shifted-component-span}
Use the marked level and two runs of
Proposition~\ref{prop:linked-comparison-rows}. On one run let
\(u=w_F+k w_L\) for an integer \(k\ge0\). Assume \(u\) has
even central degree at least \(4[E:\mathbb Q]r+4\) and
positive last weights in the respective blocks, as furnished
by the chosen \(w_F,w_L\). Let \(S_{\rm an}(u)\) be the space of analytic cusp families on the marked components, and \(\Theta_\alpha(u)\) the restrictions of the classical theta series on the common lattice family, with the prescribed scalar on each component. Then
\[
 S_{\rm an}(u)=\operatorname{span}_{\mathbb C}\{\Theta_\alpha(u)\}_\alpha.
\]

\end{proposition}

\begin{proof}
\proofstep{Apply the theta spanning theorem at the shifted weight.}
Write \(n=u.n\) for the central degree and
\(r=p_j+q_j\). The central degree of \(w_F\) is even and at least

\[
  8fr+4=4[E:\Q]r+4.
\]
 The last positive row weights of \(w_F\)
are positive. Adding \(k w_L\) preserves both conditions.
Each analytic component has the prescribed CM type, Hermitian form,
signature and polarization scalar. Use one component as reference.
Move the integral period lattices with their markings into the
self-dual lattice family of the reference space. Write \(g_\gamma\) for
the move from component \(\gamma\), \(L_\gamma\) for its moved lattice,
and \(\Gamma_D\) for the level group. The level data include the
principal polarization, the split odd level, the marked character
tuple, and the auxiliary type prime distinct from three. There are
two properties of this marked cover that matter for spanning.
Writing \([L]_{\Gamma_D}\) for the orbit of a lattice, the marking distinguishes the components by
\[
  [L_\gamma]_{\Gamma_D}=[L_{\gamma'}]_{\Gamma_D}\quad\Longrightarrow\quad \gamma=\gamma'.
\]

On the
stabilizer of \(L_\gamma\) the residual vacuum multiplier is one, and the
group transported back to the component \(\gamma\) acts trivially on
its section. Here the multiplier is what remains after the finite
character of the reduction is taken out of the Weil action.
We explain the use of these properties before giving the comparison
of the markings.

\proofstep{Comparison of theta functions in a common frame.}
Let \(P_\gamma\) denote the function representing the analytic cusp family in the chosen frame on component \(\gamma\). Move it
to the reference frame by \(g_\gamma\), dividing by the nonzero component
normalization \(C_\gamma\) in the theta function. Extend this function to
the entire orbit of \(L_\gamma\) by the inverse change of frame and the
vacuum multiplier in the automorphy law. Indeed two moves \(g_1,g_2\) leading to the same lattice differ by a stabilizer element, since
\[
  g_1L_\gamma=g_2L_\gamma\quad\Longrightarrow\quad g_2^{-1}g_1\in\operatorname{Stab}(L_\gamma).
\]
 The stabilizer properties and the cocycle law give
the same extension. Set the function to zero on every remaining
orbit.

The resulting function is bounded in the invariant norm since the multipliers
are unitary and there are finitely many chosen components.
Section~\ref{subsec:theta-row-rank} now applies to this one family.
Its inverse overlap operator and reproducing kernel give a finite
sum of theta series on the lattice family with coefficients in the
harmonic highest-weight polynomial spaces. Restricting the sum to
the \(L_\gamma\) gives all the component readings with one list of
lattices and polynomials.
\label{cor:slice-shifted-theta-1}

\end{proof}

\paragraph{The marked cover and the free lower rows.}
Write \(\mathfrak M\) for the conjugate of the marked half level
and \(\mathfrak J\) for the auxiliary type ideal. The level used
for theta spanning is 
\[
  \mathfrak m=\mathfrak M\mathfrak J,
\]
 with
\((\mathfrak M,\mathfrak J)=1\) and
\(\mathfrak J\overline{\mathfrak J}=(p_d)\).
In the group of invertible \(r\)-frames modulo \(\mathfrak m\)
take the kernel of reduction modulo \(\mathfrak M\). Chinese
remaindering identifies it with the general linear group on the
auxiliary factor. The finite test is supported where the upper
\(r\) rows reduce to the prescribed frame orbit. Its value there
is the auxiliary character in that frame.

The character is
trivial on the special linear subgroup. Indeed the finite ring
is a product of local rings whose residue fields have more than
three elements. Row reduction generates the special linear group
by elementary matrices, and conjugation by a diagonal unit whose
difference from one is a unit expresses each elementary matrix
as a commutator. In rank one the determinant is the identity.

\begin{proposition}[Theta evaluations at the marked level]
\label{prop:marked-theta-family}
For the prescribed half-level marking, the free lower rows at
\(\mathfrak m=\mathfrak M\mathfrak J\) realize the component evaluations at the shifted weights of Paragraph~\ref{par:shifted-theta}.
On component \(k\) write \(a_k(T),b_k(T)\) for the two complex evaluations of the rule \(T\) in the respective component coordinates. They are proportional by one factor on the component,
\[
 \exists C_k\in\mathbb C^\times\quad\forall T,\qquad a_k(T)=C_k b_k(T).
\]
This also applies when the desired family has different sections on its components.
\end{proposition}

\begin{proof}
\proofstep{The perfect pairing at the marked level.}
The principal polarization gives a perfect level pairing, with
the chosen primitive root of unity specifying its values. Write \(\delta_k\) for the dual half marking on component
\(k\), and choose generators \(b_i\) modulo \(\mathfrak M\)
on the reference module. At a split place a matrix \(A\)
on one half of a level pairing determines the matrix on the
opposite half by
\[
                  A^{\mathsf T}B=1,\qquad B=A^{-\mathsf T}.
\]
Thus one chooses a local basis on one half and the dual basis
on the other. Complete these frames on the auxiliary factor.
The resulting open level conditions can be satisfied by one
rational unitary move.

For this approximation use the Cayley transform
\(C(X)=(1-X)(1+X)^{-1}\) on the open set of skew-Hermitian
matrices with \(1+X\) invertible. One checks from \(X^*=-X\)
that 
\[
  C(X)^*C(X)=1.
\]
 This is the open chart of unitary
matrices with \(1+C(X)\) invertible. Its intersection with
each of the prescribed local open sets is dense by the chart
and its translates. Approximate the local parameters
simultaneously in the rational vector space of skew-Hermitian
matrices. The resulting rational unitary move satisfies the
same open level conditions.
The resulting moves take the period lattices into the indicated
genus with the same prescribed residues of \(p_d\delta_k(b_i)\)
modulo \(\mathfrak M\).

If \(gL_k=L_{k'}\) with \(g\) in the indicated level
group, the marking congruences read
\[
        p_d\,g\delta_k(b_i)
                \equiv p_d\,\delta_{k'}(b_i)\pmod{\mathfrak M}
                       \qquad(1\le i\le r).
\]
Cancel \(p_d\) and use the generating property of the \(b_i\).
The two dual markings agree on the entire module modulo
\(\mathfrak M\). The perfect pairing recovers the original
markings modulo the full rational level. Conjugating back by
the two moves gives an isomorphism of the marked polarized
period data. The uniformization identifies their moduli components.
This proves the orbit separation used above.

For a stabilizing move the same comparison puts its pullback in
the marked component group. Its determinant is an integral
norm-one unit and hence a root of unity. The principal congruence
of rational level at least three forces this torsion determinant
to be one.
Writing \(N\ge3\) for the rational principal level, the stabilizer element \(g\) satisfies
\[
  \det g\in\mu(E),\quad \det g\equiv1\pmod{N}\quad\Longrightarrow\quad\det g=1 .
\]
 For the signed vacuum, choose a symplectic basis
adapted to the split lattice and its half level.

In these
coordinates the off-diagonal blocks of the period change vanish
modulo eight and the determinant of its lower right block is
one modulo eight. These congruences are the adapted-basis
comparison for a principal-level isometry of determinant one.
For the lower right block \(D\), the comparison gives
\[
       \det D\equiv1\pmod8,\qquad \chi_4(\det D)=1 .
\]
In the squared half-characteristic transformation the
Gaussian factor is exactly the signed-vacuum coordinate change.
After taking the \(n/2\)-th power, the resulting determinant
factor cancels the one in the frame change. Here \(n\) is even. Evaluate the transformation
at a point where the theta function is nonzero. The remaining
vacuum scalar is one.
\label{cor:slice-marked-theta-cover-1}

\proofstep{Stable-range and full-level comparisons.}
There is one further level comparison in applying the stable
expansion. Its test requires zero lower rows modulo
\(\mathfrak m\), whereas the arithmetic rule tests lower rows
only modulo \(\mathfrak M\). Choose \(\alpha\) with

\[
  (\alpha)=\mathfrak J\mathfrak a
\]
 where \(\mathfrak a\) avoids
the levels and the bad primes. Multiply the lower rows by
\(\overline\alpha/\alpha\) and leave the upper rows fixed.
This is a unitary column change in the polynomial coordinates.
Let \(A\) be the transformed source lattice and \(S=\mathcal O_E^n\).
Take an integer \(M\) divisible by twice the bad and level
moduli. The construction gives \(t\) prime to \(M\) with
\(t\mathfrak J A\subseteq S\) and \(t\overline{\mathfrak J}S
\subseteq A\). Patch the lattice as
\[
              \Lambda'=\{v: Mv\in A,\ tv\in S\}.
\]
At the primes of \(M\) it equals \(S\), and at the other primes
it equals \(A\), by the two inclusions. Thus it is self-dual
and has an admissible level. On a target with shifted lattice
\(L^\sharp\), expand the moved vacuum into translates for
\(\Lambda'\). Its translation classes are indexed by

\[
  (A\otimes L^\sharp)/[(A\otimes L^\sharp)\cap
(\Lambda'\otimes L^\sharp)].
\]
 Choose representatives by
Chinese remaindering with upper corrections zero modulo
\(\mathfrak m\) and lower corrections zero modulo
\(\mathfrak M\). The upper rows still select the original
character. Conversely, keep the upper rows of any admissible
free-row argument and set the other rows to zero. This determines
its strict-test class. Subtract its moved representative from
the argument. The difference is in \(A\otimes L^\sharp\) and
determines the translation class uniquely. The even representatives
give trivial composition phase. This proves equality of the
two finite sums up to vacuum normalization. The positive factor
from projection is

\[
  [A:A\cap\Lambda']^r[\Lambda':A\cap\Lambda']^{-r/2}.
\]
The column change also multiplies the normalized vacuum by
a scalar of modulus one. Both factors are common to the targets,
so the harmonic polynomial, transformed by the same unitary
change, absorbs their ratio.
\label{cor:slice-marked-theta-cover-2}

\proofstep{A common scalar relating the two theta normalizations on each component.}
Finally, on a unit frame of component \(k\), the two normalized
vacua at the source of a short rule differ by one scalar. The
finite translate and derivative formulas therefore agree up to
this scalar for every rule simultaneously, after removing
\(u_{\Lambda,N_\Lambda}\) from each arithmetic reading. Its
modulus is fixed on that component. Write \(A_R\) for the normalized arithmetic reading of
a rule \(R\) after removal of \(u_{\Lambda,N_\Lambda}\), and
\(T_R\) for its classical reading. The vacuum comparison gives

\[
  A_R T_S=A_S T_R
\]
 for any two rules, in their common flag
character. If some \(T_S\) is nonzero, work first on its
nonvanishing locus. There \(A_S/T_S\) is holomorphic with fixed
positive modulus, so it equals a constant \(C_k\) on an open
set. The identity theorem on the connected frame domain gives
\(A_S=C_kT_S\) everywhere. For each \(R\) the cross-product identity now gives
\[
 (A_R-C_kT_R)T_S=A_R T_S-A_S T_R=0 .
\]
Cancel \(T_S\) on its nonvanishing locus; the identity theorem extends
the resulting equality over the connected frame domain. If all classical readings vanish, the vacuum
comparison makes the arithmetic readings zero too and any
scalar of the prescribed modulus serves. This is the common
\(C_k\) used in extending an arbitrary component family.
\label{cor:slice-marked-theta-cover-3}
\end{proof}


\begin{proposition}[Rational sections at the shifted weights]\label{prop:shifted-rational-cusp-span}
Let \(\mathcal S(u)\) be the rational short sections and \(K\) the coefficient field of \(C(u)\). For each \(k\ge0\),
\[
 u=w_F+k w_L,\qquad \operatorname{span}_K\mathcal S(u)=C(u).
\]
Their analytic realizations span all analytic cusp families on the run.
\end{proposition}

\begin{proof}
In the resulting classical series write \(\Lambda\) for a source
lattice and \(N_\Lambda\) for an integer exhibiting its admissible
level. The integral harmonic polynomials at this level span its
complex highest-weight space. A concrete spanning set is constructed
below. The classical evaluation of a short rule attached to an
integral polynomial \(P\), on component \(\gamma\), is the corresponding
classical theta reading times \(C_\gamma u_{\Lambda,N_\Lambda}\), with
\(u_{\Lambda,N_\Lambda}\ne0\) independent of the component \(\gamma\). Write \(\Theta_{\Lambda,\gamma}(P)\) for the classical theta
reading on component \(\gamma\), linear in the harmonic polynomial.
The common classical expansion has coefficients \(P_\Lambda\) with

\[
  C_\gamma^{-1}P_\gamma=\sum_\Lambda\Theta_{\Lambda,\gamma}(P_\Lambda).
\]
Replace each polynomial by \(P_\Lambda/u_{\Lambda,N_\Lambda}\)
and express the result in the spanning integral polynomials.
The corresponding short rules then have total reading \(P_\gamma\)
on every component, using the same coefficients on them all. Both sides can be
extended by zero off the frame domains. For an empty component set
the assertion holds with an empty sum. The transformation law
gives rational sections and the last-positive-weight condition
gives cuspidality at every good embedding. Faithfulness after
extension of \(K\) follows from the independent evaluation rows of
Section~\ref{subsec:selected-cusp-reading}. In particular the rational
span of the short rules is \(C(u)\), since a proper rational
subspace stays proper after extending scalars.
\label{cor:slice-shifted-theta-2}

\end{proof}

\paragraph{Descent of the low-mode space.}\label{par:shifted-good}
Suppose \(j<m\). For each admissible CM arc over \(K'/K\), use the
coefficient functional \(\phi:C(\kappa)\to K'\) of the specified
frame monomial and order. The arc calculation at the half-trace
cutoff \(c_1\) identifies the intersection of their kernels over
\(K\) with \(H_g\). For a complex linear combination with all low
modes zero, the coefficient equation holds at \emph{each} extension
\(\iota:K'\to\C\) of \(\sigma_0\). Indeed at that embedding the
conjugate arc realizes the same Fourier--Jacobi test on the unfolded
components, with the same cutoff. Extracting the coefficient from
the sum commutes with its finite complex linear combination.

\begin{proposition}[The rational good space and its complex extension]
\label{prop:shifted-low-mode-descent}
At \(\kappa\) the complex families with the prescribed low modes
zero are exactly the scalar extension of the rational kernel \(H_g\)
of the coefficient tests. Thus if \(S_{\rm low=0}\) denotes the complex subspace defined by those coefficient tests,
\[
 S_{\rm low=0}=H_g\otimes\mathbb C,\qquad H_g=C(\kappa)\quad\text{at the definite lift}.
\]

The rational kernel spans the same good families used in the
slice-rank and division arguments.
\end{proposition}

\begin{proof}
\proofstep{Descent by annihilation of the low modes.}
To see why these equations suffice, take a full-span expansion
of the analytic family in rational forms and choose a
\(\sigma_0(K)\)-basis \(b_1,\ldots,b_h\) of the span of its complex
coefficients. Regroup the rational forms into \(F_1,\ldots,F_h\)
so the expansion has coefficients \(b_i\). For each arc,
\(\sum_i b_i\iota(\phi(F_i))=0\) for every \(\iota\) above
\(\sigma_0\).

For any \(a\in K'\), multiply by \(\iota(a)\) and
sum over these embeddings. The result is
\[
 \sum_i b_i\,\sigma_0\!\left(
       \operatorname{Tr}_{K'/K}(a\phi(F_i))\right)=0 .
\]
Independence and the nondegenerate trace pairing give
\(\phi(F_i)=0\) for every \(i\). The same \(F_i\) work for every
arc, so all belong to \(H_g\). Conversely the readings of \(H_g\)
have zero low modes by the detector. At \(j=m\) the Hermitian
form is definite, there is no isotropic line, and \(H_g=C(\kappa)\).
Thus in both cases the interpolation, rank and division arguments
use precisely the scalar extension of the rational good space.
\label{cor:slice-shifted-theta-3}

\end{proof}

\paragraph{Small representatives of the three rational spans.}\label{par:shifted-small}
Write \(L=x^\epsilon\), \(z=frL\), \(h=frn\), and \(d\) for the
total degree of the harmonic polynomial at weight \(u\). Let \(A\)
be the sum of all the flag entries of \(u\), so that \(d\le A\) and
\(A+h+1\le |u|\). We use the selected size bounds
\[
 L\ge2,\qquad \log f,\ \log\operatorname{rd}(E),\
            \log(1+A)\le2L,\qquad \log n\le L,\qquad r\le n .
\]
They hold at each of the three weights separately by coordinatewise
domination by \(\kappa\). We give estimates independent of the
lattice class and the embedding of \(K\).

\begin{proposition}[Individual height bounds at the three weights]
\label{prop:three-shifted-spanning-sets}
Each of the three specified weights has a rational spanning set
in its scaled theta lattice in which each chosen form \(g\) satisfies
\[
 H(s_uM_0^{frn}g)\le G_c|u|.
\]
Here \(H\) uses the normalization of the proof. The size bound is independent of the
source lattice class and the coefficient embedding and is retained by
any subset chosen as a basis.
\end{proposition}

\begin{proof}
On a unimodular source lattice \(\Lambda\) of rank \(n\), choose
\(2fn\) rationally independent vectors whose norm at every embedding
is at most \(\rho=2f^{3/2}n^2\operatorname{rd}(E)\).
Indeed for \(v\ne0\) in its dual for the trace of the Hermitian
pairing, \(\langle v,\Lambda\rangle\subseteq\mathfrak D_E^{-1}\).
Unimodularity then implies \(\mathfrak D_E v\subseteq\Lambda\).
Hence \(\langle v,v\rangle\in\mathfrak D_E^{-2}\), using invariance
of the different under conjugation. Its positive conjugates have
product at least \(|\operatorname{disc}E|^{-2}\). Their sum is
at least \(2f/\operatorname{rd}(E)^2\).

For the \(2f\) positive conjugates \(a_\sigma\), the arithmetic--geometric mean bound gives
\[
  \sum_\sigma a_\sigma\ge 2f\Big(\prod_\sigma a_\sigma\Big)^{1/(2f)}\ge \frac{2f}{\operatorname{rd}(E)^2}.
\]
 Euclidean transference in
dimension \(2fn\) bounds the last successive minimum of the
original lattice by \(2fn\,\operatorname{rd}(E)/\sqrt{2f}\)
in the trace norm, which suffices for the displayed \(\rho\).
The source lattices are full and finitely generated since the
admissible-level conditions sandwich them between two scalar
multiples of the standard lattice.
\cite{banaszczyk}\label{cor:slice-short-heights-1}

\proofstep{The bounded harmonic grid.}
Use these vectors for the two sets of row letters in the harmonic
polynomial of the short rule. Here is a finite interpolation argument
which controls their size even for very large row weights. Write
\(a+b=r\) for the two block sizes. Start with three lists of lengths
\(a,a,b\) of vectors \(e_i,e'_i,g_j\). Let 
\[
  B_{li}=\langle
e'_l,e_i\rangle,
\]
 with the pairing linear in the first variable.

In the second block use the letters
\[
 f_j=(\det B)g_j-\sum_{i,l}\langle g_j,e_i\rangle
                         (\operatorname{adj}B)_{il}e'_l .
\]
They are integral and \(\langle f_j,e_i\rangle=0\) identically,
including for singular \(B\). Use the products of leading minors
with the row-gap exponents, on \(e\) and the conjugates of \(f\).
They satisfy the flag weight and are harmonic because each mixed
contraction pairs orthogonal letters. In the opposite convention
conjugate both the lattice and its letters before making the
construction. This gives the integral polynomial in the symmetric
algebra of the original lattice.

\proofstep{Spanning after extension to the splitting field.}
For spanning, over a splitting field the independent embeddings
and their conjugates are the separate factors of

\[
  E\otimes_{\Q}\overline{\Q}.
\]
 The simultaneous orbit of a highest
minor product spans its highest-weight harmonic space under column
changes: on each factor it generates the irreducible trace-free
Schur representation with the two prescribed partitions. The
condition \(a+b\le n\) permits the two opposite blocks of an
invertible column frame and its dual. Such orthogonal pairs on the
open orbit are obtained in the displayed parametrization by choosing
the primed list dual to the first block and setting the last list
equal to the desired perpendicular vectors. This prescription can
be made in the split factors because all the signed embeddings
are distinct.

Express the three lists in the rational basis of
\(2fn\) short vectors. The map from coefficients to the minor
product is polynomial of degree at most 
\[
  D=(2r+1)d:
\]
the adjugate expression has degree at most \(2a+1\). Any linear
functional annihilating its values at coordinates in
\(\{0,\ldots,D\}\) annihilates the polynomial, by successive
univariate interpolation. The grid therefore spans the same
space in characteristic zero. Degree zero uses the constant
polynomial when present, and the zero space uses the empty set.

With \(d>0\), each grid vector has length at most

\[
  R=2fnD\rho
\]
 at every embedding. Cauchy--Schwarz and determinant
expansion bound the length of every projected letter by
\((a+1)!R^{2a+1}\). Each linear form on \(n\) coordinates has
coefficient sum at most \(\sqrt n\) times its Euclidean norm.
The determinant expansion in a leading minor of size \(i\) costs
at most \(i!\le r^i\). Submultiplicativity thus gives, at every
embedding,
\[
 \|P\|_1\le\{r\sqrt n\,(r+1)!\,R^{2r+1}\}^{d},
 \qquad \log\|P\|_1\le140rLd .
\]
Indeed \(R\le e^{17L}\), \((r+1)!\le e^{4rL}\), and
\(r\sqrt n\le e^{2L}\) under the size bounds above. The rule is
linear in \(P\) for fixed lattice, marking and translates by the
covariant derivative expression on each chart. Replacing its
polynomial by this grid family preserves the rational span.
\label{cor:slice-short-heights-2}

\proofstep{The passage from coefficient height to Petersson norm.}
Here are details of the passage from coefficient size to the
supremum at an embedding. For a nonzero section choose an
orthonormal frame at more than half the supremum of its invariant
norm before the Lie algebra normalization. At that frame the first
minimum \(\lambda\) of the period lattice satisfies
\(\log(1/\lambda^2)\le80L\).
In the rank-one period argument of the integral lower norm, take
spread bound \(w\), transvection period \(M\), reference-lattice
clearing integer \(N_A\), and polarization factor \(k_0\le2\).
The period bound is
\[
 \lambda^2\ \ge\ k_0 e_0/N_A^2,\qquad
 e_0=\min\{1,\pi/[M(fw+2)k_0\log 2]\}.
\]

In fact \(N_A v\) belongs to the reference lattice and if every
majorant of that vector were below \(e_0\), the negative imaginary
shift by \(1/(2k_0e_0)\) and Schwarz's lemma at the vanishing cusp
would contradict the half-supremum choice. In the definite case a
nonzero isotropic vector is already excluded. The rank-one data of
the run give \[
 \log M\le30L,\ \log(fw+2)\le5L,\ \log N_A\le4L,\qquad \log(1/k_0)\le3L.
\] Thus \(-\log e_0\le36L\), which gives
the asserted bound. Conjugate embeddings use the conjugate
uniformization with the same norm.
\label{cor:slice-short-heights-3}

In normalized Gaussian coordinates the tensor period lattice has
minimum at least \(\lambda/\rho\). A nonzero tensor has nonzero
contraction with one of the \(2fn\) short vectors, and that contraction
is a period of length at most \(\rho\) times the tensor norm. Its
torus is principally polarized and has covolume one in these
coordinates. For a theta function of mean squared norm one in
dimension \(h\), the polydisc of radius
\[
 b_0=\min\{1/\sqrt{\pi h},\lambda/(4\rho\sqrt h)\}
\]
in every complex coordinate injects into the torus. Translate the
function by the unitary Heisenberg operator to put its evaluation
at zero. The submean bound for its squared modulus uses the
polydisc volume \((\pi b_0^2)^h\). The Gaussian density changes
on that polydisc by at most \(e\).

Taking square roots in the submean estimate, the invariant supremum of the normalized function \(h_0\) satisfies
\[
  \|h_0\|_{\sup}\le\left(\frac{e}{(\pi b_0^2)^h}\right)^{1/2}=e^{1/2}(\sqrt\pi\,b_0)^{-h}.
\]
The logarithm of the last factor is at most \(55hL\), by the
bounds on \(\lambda,\rho,f,r,n\). For a derivative of multi-index
\(\alpha\) of degree \(d\), Cauchy's formula in the active variables
with radii \(\sqrt{\alpha_i/\pi}\) costs at most
\((\pi e d)^{d/2}\) times that supremum, since
\(\alpha_i!\le\alpha_i^{\alpha_i}\). A homogeneous polynomial costs
its coefficient sum times this factor. Changing from normalized
letters to the original letters multiplies this by \(t^{d/2}\),
where \(t=2\xi_R\) is the polarization scale. The Lie norm of the
run contributes \((\pi t)^{-d/2}\).
The factors involving the polarization scale \(t\) cancel, leaving
\[
  t^{d/2}(\pi t)^{-d/2}=\pi^{-d/2}.
\]

\proofstep{The unit-mean normalization.}
For clarity, the unit mean normalization has no class-index cost in
the short reading. Write \(J=[\Lambda:\Lambda\cap\mathcal O_E^n]^r\)
with the index taken on the side specified by the tensor translate
count. Equivalently it is the number of cosets in the projected
vacuum sum before averaging. Distinct cosets have orthogonal
translate characters on a common finite cover, so their sum has
squared mean \(J\) times that of the tensor power of the initial
vacuum. Projection divides the sum by \(J\). After removal of the
initial vacuum metric its norm consequently has the extra factor
\(1/\sqrt J\). In the rule the torsion count divided by its initial
level index equals \(J\), whereas the arithmetic normalizing
factor \(\beta_\Lambda\) has \(|\sigma\beta_\Lambda|^2=J\) and
occurs inversely. The resulting factor is

\[
  J/(\sqrt J\sqrt J)=1.
\]
 The residual finite-character scalar
has modulus one. The finite mean of translates can be bounded
term by term, since all its coefficients have modulus at most
\(p_{\rm type}^{\,h}\), with \(\log p_{\rm type}\le L\).
These are the two normalizations in the short rule on the
marked level.
\label{cor:slice-short-heights-4}

Including the factor two from evaluation near the supremum, the
preceding estimates give each unscaled short representative the
bound
\[
 \log\|g\|_{\sup,\sigma}
 \le 140rLA+10LA+200hL\le400z|u|.
\]
In particular no count of the potentially large interpolation
grid occurs in this bound. Finite dimension selects a finite spanning
set \(\operatorname{Sh}(u)\) of the separate short representatives.

\proofstep{Clearing the product lattice.}
To put these forms in the lattice used for products, let
\(M_0=g_0g_1\), and let \(s_u\) be the selected nonzero weight
scalar, namely \(s_F,s_0,s_\kappa\) at the respective weights.
The short rule contributes \(M_0^{frn}g\) to the unscaled
theta lattice used here. Thus the actual integral columns
are \(s_uM_0^{frn}g\). Their rational span is unchanged.

The common scale bounds \(\log M_0\le\lambda_*\) and
\(\log|\sigma s_u|\le2(\lambda_*+2\log p)|u|\), where \(p\)
is the run prime, give the mean log bound
\[
 H(s_uM_0^{frn}g)
 \le [400z+2(\lambda_*+2\log p)+\lambda_*]\,|u|
 \le G_c |u|.
\]
For the last inequality the growth allowance contains \(1009z\)
and \(2fr(\lambda_*+2\log p)\), as well as the nonnegative
coefficient-field and lattice terms. Here \(fr\ge 12\).
Every subset selected by independence at a later stage inherits
the same individual bound.
\label{cor:slice-short-heights-5}
\end{proof}



\paragraph{Descent of the containment.}

The three spans have now been defined over the coefficient field. The containment needed in the index estimate concerns their saturated intersections with the integral row.

\begin{proposition}[The saturated containment on the selected row]
\label{prop:descent-of-slice-division}
With the rational good space of Proposition~\ref{prop:shifted-low-mode-descent},
the containment provided by Proposition~\ref{prop:lift-and-divide-at-slice}
holds over the generic coefficient field \(K(X)\). Writing \(M_g\) for the integral good submodule, and taking saturation inside the completed cusp lattice, one has
\[
 \operatorname{Sat}(M_g)\subseteq\operatorname{Sat}(N_1+Z).
\]

Finite indices inside these saturations require separate bounds.
\end{proposition}

\begin{proof}
A rational form of the selected good space has the required
low modes zero at \(\kappa\) by the maximal cusp comparison.
The preceding division writes its reading in the complex
span of the leading products and the boundary products.
Work first over \(K(X)\), the coefficient field with the
cutting coefficients indeterminate. The products are
polynomials in \(X\) in the rational cusp space.

At the
chosen algebraically independent complex tuple, the
full reading on that space is faithful after scalar
extension. Thus the ranks of the columns with and
without the rational form are exactly the ranks of the
original matrices over \(K(X)\), as polynomial minors
vanish at this tuple precisely when they are zero
polynomials.
Writing \(A\) for the column matrix and \(v\) for the rational form to be tested, membership is the minor test
\[
  v\in\operatorname{im}A\quad\Longleftrightarrow\quad \operatorname{rank}[A\mid v]=\operatorname{rank}A .
\]
 This gives the containment over \(K(X)\).
Extend scalars to the fraction field of the Gauss DVR
at each finite place. Intersecting with the completed cusp
lattice then gives containment of the saturated good
lattice in the saturation of \(N_1+Z\). The indices inside
those saturations still require the finite-place comparisons.
\label{cor:slice-short-division-3}
\end{proof}
